% !TeX root = ../../../main.tex
% chapters/sections/chapter5/subsections/two_plane/cs.tex

\subsection{CS（消費供給）曲線}
\label{sec:chapter5_two_plane_cs}
CS（消費供給）曲線は消費者物価指数（CPI）\(p_t^{H\to W}\)を総消費指数\(c_t^{H\to W}\)の関数として
記述したものであり、YS曲線に外部群の式を統合することで導出される。

\subsubsection{CS曲線導出の準備}
\label{sec:chapter5_two_plane_cs_preparation}
次節におけるCS曲線の導出にあたりいくつか式を準備しておく。
消費者物価指数（CPI）\(p_t^{H\to W}\)の式
\eqref{form:chapter5_two_plane_classification_external_p_H_W_eq}
に対して、外国の一物一価の法則
\eqref{form:chapter3_assumptions_lop_foreign_p_f_eq}、
名目為替レートの決定式
\eqref{appendixTODO}
\begin{equation}
\label{form:chapter5_two_plane_cs_preparation_e_slash_star_eq}
e_t^{/*}=\operatorname{E}_t[\mathcal{B}_t\Lambda_{\infty}]\lambda_t^{F/*}/\lambda_t^H
\end{equation}
および所得の限界効用の式
\eqref{form:chapter5_two_plane_classification_external_lambda_H_p_H_W_c_H_W_eq}
を順次代入する。
\begin{equation}
\label{form:chapter5_two_plane_cs_preparation_p_H_W_eq}
    \begin{aligned}
        p_t^{H\to W}
        &=(p_t^H)^{\alpha^H}(p_t^F)^{1-\alpha^H} \\
        &=(p_t^H)^{\alpha^H}(e_t^{/*}p_t^{F*})^{1-\alpha^H} \\
        &=(p_t^H)^{\alpha^H}\left(\operatorname{E}_t[\mathcal{B}_t\Lambda_{\infty}]\frac{\lambda_t^{F/*}}{\lambda_t^H}p_t^{F*}\right)^{1-\alpha^H} \\
        &=(p_t^H)^{\alpha^H}\left(\operatorname{E}_t[\mathcal{B}_t\Lambda_{\infty}]\lambda_t^{F/*}p_t^{F*}(p_t^{H\to W}c_t^{H\to W})\right)^{1-\alpha^H}
    \end{aligned}
\end{equation}
自国金融政策の影響を受けない外部変数群をまとめて
\(Z_t\equiv\operatorname{E}_t[\mathcal{B}_t\Lambda_{\infty}]\lambda_t^{F/*}p_t^{F*}\)
とおき、両辺を\((p_t^{H\to W})^{1-\alpha^H}\)で割って\(1/\alpha^H\)乗すると以下のようになる。
\begin{equation}
\label{form:chapter5_two_plane_cs_preparation_p_H_W_modified_eq}
    p_t^{H\to W}=p_t^H(Z_tc_t^{H\to W})^{\frac{1-\alpha^H}{\alpha^H}}
\end{equation}
これを対数線形化することで以下の1つ目の式が得られる。
\begin{equation}
\label{form:chapter5_two_plane_cs_preparation_p_H_W_modified2_eq}
    \hat{p}_t^{H\to W}=\hat{p}_t^H+\frac{1-\alpha^H}{\alpha^H}(\hat{Z}_t+\hat{c}_t^{H\to W})
\end{equation}
さらにこの式とこの式で\(t\)を\(t-1\)としたものの差をとることにより、
消費者物価インフレ率\(\hat{\pi}_t^{H\to W}\)と
生産者物価インフレ率\(\hat{\pi}_t^H\)の乖離を示す2つ目の式が得られる。
\begin{equation}
\label{form:chapter5_two_plane_cs_preparation_pi_H_W_eq}
    \hat{\pi}_t^{H\to W}=\hat{\pi}_t^H+\frac{1-\alpha^H}{\alpha^H}(\hat{Z}_t-\hat{Z}_{t-1}+\hat{c}_t^{H\to W}-\hat{c}_{t-1}^{H\to W})
\end{equation}
次に名目GDP\(p_t^Hy_t^H\)と名目総生産\(p_t^{H\to W}c_t^{H\to W}\)の関係式
\eqref{form:chapter5_two_plane_yd_preparation_p_H_y_H_eq}
は以下のようであった。
\begin{equation}
\label{form:chapter5_two_plane_cs_preparation_p_H_y_H_eq}
    p_t^Hy_t^H=X_tp_t^{H\to W}c_t^{H\to W}
\end{equation}
ただし\(X_t\equiv\alpha^H+(1-\alpha^F)\operatorname{E}_t[\mathcal{B}_t\Lambda_{\infty}]\)であり、
この\(X_t\)は遮断効果により自国金融政策の影響を受けない。
これを対数線形化する。
\begin{equation}
\label{form:chapter5_two_plane_cs_preparation_p_H_y_H_modified_eq}
    \hat{p}_t^H+\hat{y}_t^H=\hat{X}_t+\hat{p}_t^{H\to W}+\hat{c}_t^{H\to W}
\end{equation}
これを\(\hat{y}_t^H\)について解き、先ほど導いた1つ目の式
\eqref{form:chapter5_two_plane_cs_preparation_p_H_W_modified2_eq}
を代入すると以下の3つ目の式が得られる。
\begin{equation}
\label{form:chapter5_two_plane_cs_preparation_p_H_y_H_modified2_eq}
\begin{aligned}
    \hat{y}_t^H
    &=\hat{X}_t+(\hat{p}_t^{H\to W}-\hat{p}_t^H)+\hat{c}_t^{H\to W} \\
    &=\hat{X}_t+\frac{1-\alpha^H}{\alpha^H}(\hat{Z}_t+\hat{c}_t^{H\to W})+\hat{c}_t^{H\to W} \\
    &=\hat{X}_t+\frac{1-\alpha^H}{\alpha^H}\hat{Z}_t+\frac{1}{\alpha^H}\hat{c}_t^{H\to W}
\end{aligned}
\end{equation}

\subsubsection{CS曲線の導出}
\label{sec:chapter5_two_plane_cs_derivation}
前節で導いた3つ目の式
\eqref{form:chapter5_two_plane_cs_preparation_p_H_y_H_modified2_eq}
をYS-\(\pi\)曲線
\eqref{form:chapter5_two_plane_ys_derivation_pi_H_modified_eq}
の右辺に代入する。
\begin{equation}
\label{form:chapter5_two_plane_cs_derivation_pi_H_eq}
\begin{aligned}
    \hat{\pi}_t^H
    &=\beta_{ss}^H\operatorname{E}_t[\hat{\pi}_{t+1}^H]+\kappa^H\left(2\left(\hat{X}_t+\frac{1-\alpha^H}{\alpha^H}\hat{Z}_t+\frac{1}{\alpha^H}\hat{c}_t^{H\to W}\right)-2\hat{a}_t^H-\hat{X}_t-\hat{\tau}_t^H\right) \\
    &=\beta_{ss}^H\operatorname{E}_t[\hat{\pi}_{t+1}^H]+\kappa^H(\hat{\Omega}_t-2\hat{a}_t^H)+\frac{2\kappa^H}{\alpha^H}\hat{c}_t^{H\to W}
\end{aligned}
\end{equation}
ここで自国金融政策の影響を受けない外生変数群（\(\hat{a}_t^H\)を除く）をまとめて
\(\hat{\Omega}_t\equiv\hat{X}_t+\frac{2(1-\alpha^H)}{\alpha^H}\hat{Z}_t-\hat{\tau}_t^H\)
とおいた。
これを前節の2つ目の式
\eqref{form:chapter5_two_plane_cs_preparation_pi_H_W_eq}
の右辺に代入すると、
CS曲線の\(\hat{\pi}_t^{H\to W}\)を用いた表現である\textbf{CS-\(\hat{\pi}\)曲線}が得られる。
\begin{equation}
\label{form:chapter5_two_plane_cs_derivation_pi_H_W_modified_eq}
    \hat{\pi}_t^{H\to W}=\beta_{ss}^H\operatorname{E}_t[\hat{\pi}_{t+1}^H]+\kappa^H(\hat{\Omega}_t-2\hat{a}_t^H)+\frac{1-\alpha^H}{\alpha^H}(\hat{Z}_t-\hat{Z}_{t-1})+\left(\frac{2\kappa^H}{\alpha^H}+\frac{1-\alpha^H}{\alpha^H}\right)\hat{c}_t^{H\to W}-\frac{1-\alpha^H}{\alpha^H}\hat{c}_{t-1}^{H\to W}
\end{equation}
次に2平面分析に向けてこれを物価水準ベースの式に変換する。
YS-\(\hat{p}\)曲線
\eqref{form:chapter5_two_plane_ys_derivation_p_H_modified_eq}
に対しても同様に、前節で導いた3つ目の式
\eqref{form:chapter5_two_plane_cs_preparation_p_H_y_H_modified2_eq}
を代入して整理する。
\begin{equation}
\label{form:chapter5_two_plane_cs_derivation_p_H_eq}
    \hat{p}_t^H=\hat{p}_{t-1}^H+\beta_{ss}^H\operatorname{E}_t[\hat{\pi}_{t+1}^H]+\kappa^H(\hat{\Omega}_t-2\hat{a}_t^H)+\frac{2\kappa^H}{\alpha^H}\hat{c}_t^{H\to W}
\end{equation}
これを前節の1つ目の式
\eqref{form:chapter5_two_plane_cs_preparation_p_H_W_modified2_eq}
の右辺に代入し整理すると、
CS曲線の\(\hat{p}_t^{H\to W}\)を用いた表現である\textbf{CS-\(\hat{p}\)曲線}が導出される。
\begin{equation}
\label{form:chapter5_two_plane_cs_derivation_p_H_W_eq}
    \hat{p}_t^{H\to W}=\hat{p}_{t-1}^H+\beta_{ss}^H\operatorname{E}_t[\hat{\pi}_{t+1}^H]+\kappa^H(\hat{\Omega}_t-2\hat{a}_t^H)+\frac{1-\alpha^H}{\alpha^H}\hat{Z}_t+\left(\frac{2\kappa^H}{\alpha^H}+\frac{1-\alpha^H}{\alpha^H}\right)\hat{c}_t^{H\to W}
\end{equation}
価格の粘着性パラメータ\(\kappa^H\)や主観的割引因子の定常値\(\beta_{ss}^H\)は正の定数であり、
自国財の消費割合\(\alpha^H\)は\(0<\alpha^H<1\)を満たす。
これにより\(\hat{c}_t^{H\to W}\)の係数は正の値をとるため
消費者物価指数\(\hat{p}_t^{H\to W}\)は総消費指数\(\hat{c}_t^{H\to W}\)の増加関数であることがわかる。
したがって\((c_t^{H\to W},p_t^{H\to W})\)平面においてCS曲線は\textbf{右上がり}の曲線として描かれる。
CS-\(\hat{p}\)曲線に含まれる\(\hat{Z}_t\)、\(\hat{\Omega}_t\)および\(\hat{a}_t^H\)は自国金融政策の影響を受けず、
また\(\lambda_t^{F/*}\)および\(p_t^{F*}\)も遮断効果により自国金融政策の影響を受けない。
したがって金融政策の分析においてはこれらの項の動きは考慮しなくてよい。

\subsubsection{右上がりのCS-\(\hat{p}\)曲線と垂直なCS-\(\hat{p}\)曲線の幾何学的な関係}
\label{sec:chapter5_two_plane_cs_potential}
ここで\ref{chap:chapter5}におけるシミュレーション結果の解釈の準備として、
粘着価格\(\xi^H=定数\)による右上がりのCS-\(\hat{p}\)曲線と
完全伸縮価格\(\xi^H\to0\)による垂直なCS-\(\hat{p}\)曲線との幾何学的な関係について考察しておく。
粘着価格における右上がりのCS-\(\hat{p}\)曲線
\eqref{form:chapter5_two_plane_cs_derivation_p_H_W_eq}
を総消費指数\(\hat{c}_t^{H\to W}\)について解くと以下のようになる。
\begin{equation}
\label{form:chapter5_two_plane_cs_potential_c_H_W_eq}
    \begin{aligned}
        \hat{c}_t^{H\to W}&=-\frac{\alpha^H\kappa^H}{2\kappa^H+1-\alpha^H}(\hat{\Omega}_t-2\hat{a}_t^H)-\frac{1-\alpha^H}{2\kappa^H+1-\alpha^H}\hat{Z}_t+\frac{\alpha^H}{2\kappa^H+1-\alpha^H}(\hat{p}_t^{H\to W}-\hat{p}_{t-1}^H-\beta_{ss}^H\operatorname{E}_t[\hat{\pi}_{t+1}^H]) \\
        &=-\frac{\alpha^H}{2+\frac{1-\alpha^H}{\kappa^H}}(\hat{\Omega}_t-2\hat{a}_t^H)-\frac{\frac{1-\alpha^H}{\kappa^H}}{2+\frac{1-\alpha^H}{\kappa^H}}\hat{Z}_t+\frac{\frac{\alpha^H}{\kappa^H}}{2+\frac{1-\alpha^H}{\kappa^H}}(\hat{p}_t^{H\to W}-\hat{p}_{t-1}^H-\beta_{ss}^H\operatorname{E}_t[\hat{\pi}_{t+1}^H])
    \end{aligned}
\end{equation}
ここで価格が完全に伸縮的である場合を考え、両辺について\(\lim_{\xi^H\to0}\)をとる。
すると式
\eqref{form:chapter3_producers_stage2_home_kappa_eq}
より\(\kappa^H\to\infty\)よって右辺第2項および第3項は\(0\)となり、
第1項の係数は\(-\alpha^H/2\)に収束するため以下の垂直なCS-\(\hat{p}\)曲線が得られる。
\begin{equation}
\label{form:chapter5_two_plane_cs_potential_c_H_W_pot_eq}
    \hat{c}_t^{H\to W}=\alpha^H\hat{a}_t^H-\frac{\alpha^H}{2}\hat{\Omega}_t
\end{equation}
この垂直なCS-\(\hat{p}\)曲線によって定まる総消費指数\(\hat{c}_t^{H\to W}\)を潜在消費と呼び
\(\hat{c}_t^{H\to W,pot}\)により表す。
これら2つのCS-\(\hat{p}\)曲線はYS-\(\hat{p}\)曲線とは異なり\(\hat{c}_t^{H\to W}\)切片が
共通していないことに注意する。
遮断効果により自国金融政策および\(a_t^H\)は\(Z_t\)および\(\Omega_t\)に影響を与えないのであった。
そのため特に\(a\)ショック時には\(\hat{Z}_t=\hat{\Omega}_t=0\)となる。
しかしこのとき2つのCS-\(\hat{p}\)曲線の\(\hat{c}_t^{H\to W}\)切片はそれぞれ
\(\frac{2\alpha^H\kappa^H}{2\kappa^H+1-\alpha^H}\hat{a}_t^H\)および
\(\alpha^H\hat{a}_t^H\)となり一致しない。